\section{Optical Conveyor Belt}
\label{sec:conveyor_belt}

The optical dipole trap discussed in the previous section provides a static conservative potential in which atoms are confined near the focus of a single laser beam. Transporting the trapped cloud over macroscopic distances requires a different architecture, in which the position of the potential minima can be controlled dynamically. The "optical conveyor belt" fulfills this requirement by replacing the single beam with \CCn{two counter-propagating, phase-coherent beams}{``Phase-coherent'' da solo non basta: i due fasci devono anche essere CO-POLARIZZATI perche' si formi un reticolo di intensita'. Con polarizzazioni ortogonali l'intensita' totale e' uniforme e si ottiene un gradiente di polarizzazione, cioe' una melassa lin-perp-lin, non un conveyor belt. E' anche un requisito di allineamento reale in laboratorio.}\MDt{}{anceh qui sfori a destra, controlla i margini}\CCmd{titolo 1.3.1}{CONFERMO per il PDF di revisione, ma il contrario nel layout finale. Misurato: 22.48 pt di sfondamento con i margini larghi, ZERO in A4. E' l'opposto del titolo 1.1.1, che sfora solo in A4. Verificato anche che non e' colpa di \texttt{\textbackslash uwave}: lo sfondamento e' identico con e senza. Coperto comunque dal \texttt{\textbackslash raggedright}. Nota di metodo, CC: questa nota stava dentro l'argomento di \texttt{\textbackslash subsection}, che e' un argomento mobile. Il testo della nota finiva stampato verbatim nell'indice, hyperref segnalava ``Token not allowed in a PDF string'' e il \texttt{\textbackslash todo} veniva eseguito due volte, una dall'indice. L'ho spostata subito DOPO il titolo: il testo della nota e' il tuo, invariato; si perde solo l'evidenziazione della singola parola.}: their interference generates a standing wave whose nodes and antinodes can be set into motion by imposing a controlled frequency difference between the two arms. This section derives the geometry of the standing-wave potential, the kinematics of the moving lattice, and the acceleration limits that bound adiabatic transport.

\subsection{Interference of Counter-Propagating Gaussian Beams and the Stationary Lattice}

The fundamental mechanism of the optical conveyor belt relies on the periodic intensity pattern generated by the interference of two counter-propagating laser beams \cite{Kuhr:03}. The resulting standing wave acts as a one-dimensional optical lattice, imposing a tight spatial confinement on atoms via the optical dipole force.

Assuming two coherent Gaussian beams in the fundamental transverse mode propagating along the $z$-axis, the respective electric fields can be modeled as follows. The first beam, propagating towards $+z$ with angular frequency $\omega_1$, is described by:
\begin{equation}
E_1(z, \rho, t) = E_0 \frac{w_0}{w(z)} \exp\left(-\frac{\rho^2}{w^2(z)}\right) \exp\left(i(k z - \omega_1 t + \phi(z))\right)
\label{eq:beam1}
\end{equation}
where $E_0$ denotes the peak field amplitude, $k = 2\pi/\lambda$ the optical wavenumber, \CCn{$\phi(z)$ the Gouy phase shift}{Due cose. Il valore di $\phi(z)$ non viene mai dato: e' $\arctan(z/z_R)$. E soprattutto, nelle espressioni dei due fasci non compare alcun termine di curvatura del fronte d'onda, $\exp(ik\rho^{2}/2R(z))$: piu' avanti si legge ``dropping wavefront curvature and Gouy phase'', ma la curvatura non c'era da eliminare. O la si aggiunge, o si corregge quella frase.}, and $\rho=\sqrt{x^2+y^2}$ the radial coordinate. The beam radius $w(z)$ evolves along the propagation axis according to:
\begin{equation}
w^2(z) = w_0^2 \left(1 + \frac{z^2}{z_R^2}\right),
\label{eq:beam_waist}
\end{equation}
with $w_0$ the waist at the focal plane and $z_R = \pi w_0^2/\lambda$ the Rayleigh length. The second field, counter-propagating towards $-z$ at frequency $\omega_2$, is given by:
\begin{equation}
E_2(z, \rho, t) = E_0 \frac{w_0}{w(z)} \exp\left(-\frac{\rho^2}{w^2(z)}\right) \exp\left(i(-k z - \omega_2 t - \phi(z))\right).
\label{eq:beam2}
\end{equation}

To determine the local intensity profile for degenerate frequencies ($\omega_1 = \omega_2 = \omega$), the coherent superposition $E_{tot} = E_1 + E_2$ is evaluated. Restricting the analysis to the spatial region localized around the focal point ($z \ll z_R$), the variation of the beam waist and the Gouy phase shift become negligible. The superimposed electric field simplifies to:
\begin{equation}
E_{tot}(z, \rho, t) \approx E_0 \frac{w_0}{w(z)} \exp\left(-\frac{\rho^2}{w^2(z)}\right) e^{-i\omega t} \left[ e^{ikz} + e^{-ikz} \right].
\end{equation}
Applying Euler's relation, the spatial interference manifests as a real cosine envelope $2\cos(kz)$. The effective optical intensity $I(z, \rho)$ experienced by the atoms is proportional to the time-averaged squared amplitude of the field, $\langle |E_{tot}|^2 \rangle$, which yields a stationary standing wave:
\begin{equation}
I(z,\rho) = I_{\max} \frac{w_0^2}{w^2(z)} \exp\left(-\frac{2\rho^2}{w^2(z)}\right) \cos^2(kz).
\label{eq:stationary_intensity}
\end{equation}
The peak intensity $I_{\max} = 4P/(\pi w_0^2)$ follows from the total optical power $P$ distributed across the two beams. \CCn{Assuming each beam carries power}{Il ``Assuming'' aggiunto risolve meta' del rilievo: adesso l'ipotesi e' dichiarata invece che implicita. Resta la conseguenza, che vale una riga: se i due bracci non hanno la stessa potenza l'intensita' e' $I_1+I_2+2\sqrt{I_1I_2}\cos(2kz)$ e il contrasto del reticolo scende a $2\sqrt{I_1I_2}/(I_1+I_2)<1$, quindi i minimi non vanno a zero e la profondita' utile cala. Nel setup reale il braccio di ritorno e' quasi sempre piu' debole, quindi non e' un caso di scuola.} $P/2$ and has a peak intensity $2(P/2)/(\pi w_0^2) = P/(\pi w_0^2)$; the additional factor of $4$ in $I_{\max}$ originates from the constructive interference at the antinodes, where the two field amplitudes add in phase.

In the far-detuned regime, where the laser frequency is significantly shifted from any atomic resonance, the resulting dipole force is conservative. The optical potential energy $U(\rho,z)$ is therefore directly proportional to the local intensity distribution:
\begin{equation}
U(\rho,z) = U_0 \frac{w_0^2}{w^2(z)} \exp\left(-\frac{2\rho^2}{w^2(z)}\right) \cos^2(kz),
\label{eq:stationary_potential}
\end{equation}
where $U_0$ defines the maximum depth of the individual potential wells, with $U_0 < 0$ for a red-detuned, attractive trap.

The harmonic spatial modulation introduced by the $\cos^2(kz)$ term imposes a longitudinal periodicity, organizing the optical potential into a one-dimensional lattice of trapping sites with a lattice constant of $\lambda/2$.

Within a potential well, at low temperature, the atomic motion is restricted to a small region around the trap minimum at $z=0$ and $\rho=0$, allowing the harmonic approximation. Applying a second-order Taylor expansion to Eq.~\eqref{eq:stationary_potential} using $\cos^2(kz) \approx 1 - (kz)^2$ and $\exp(-2\rho^2/w_0^2) \approx 1 - 2\rho^2/w_0^2$, the potential takes the quadratic form
\begin{equation}
U(\rho,z) \approx U_0 \left(1 - \frac{2\rho^2}{w_0^2}\right)(1 - k^2 z^2) \approx U_0 - U_0 k^2 z^2 - U_0 \frac{2\rho^2}{w_0^2}.
\end{equation}
Equating these leading-order restoring terms to the classical energies of a harmonic oscillator, $\tfrac{1}{2}\kappa_z z^2$ and $\tfrac{1}{2}\kappa_{rad} \rho^2$, yields the effective spring constants $\kappa_z = 2|U_0|k^2$ and $\kappa_{rad} = 4|U_0|/w_0^2$. The characteristic axial and radial oscillation frequencies therefore read:
\begin{equation}
\Omega_z = \sqrt{\frac{2 |U_0| k^2}{m}}, \qquad \Omega_{rad} = \sqrt{\frac{4 |U_0|}{m w_0^2}}.
\label{eq:oscillation_frequencies}
\end{equation}
These characteristic frequencies establish the dynamical timescales of the system and set the bound within which external parameters can be varied while preserving adiabaticity.

\subsection{Geometry of the Trapping Potential}

\begin{wrapfigure}{R}{0.25\textwidth}
\centering
\includegraphics[width=0.20\textwidth]{Cap_1_figures/pancake.png}
\caption{The strongly anisotropic, disc-shaped potential wells of the one-dimensional optical lattice.}
\label{fig:pancake}
\end{wrapfigure}

The ratio of the two oscillation frequencies exhibits a geometric dependence, independent of the overall trap depth $U_0$:
\begin{equation}
\frac{\Omega_z}{\Omega_{rad}} = \frac{k w_0}{\sqrt{2}} = \frac{\sqrt{2}\,\pi w_0}{\lambda}.
\end{equation}
In typical experimental configurations the beam waist exceeds the optical wavelength ($w_0 \gg \lambda$), so that the longitudinal confinement dominates over the transverse confinement ($\Omega_z \gg \Omega_{rad}$). The standing wave then imposes quasi-two-dimensional potential wells with a strongly anisotropic, disc-shaped geometry: along the $z$-axis the atomic motion is localized to a fraction of $\lambda$, while the transverse motion retains a comparatively broad spatial extent. This anisotropy (provided that the potential is sufficiently deep) suppresses tunneling between adjacent lattice sites while leaving the transverse dynamics only weakly confined.
A schematic representation of the resulting disc-shaped potential wells is shown in Fig.~\ref{fig:pancake}.

\subsection{Kinematics of the Moving Standing Wave}

Translating the stationary lattice into an active transport potential requires inducing a macroscopic translation of the interference pattern. This is achieved by imposing a controlled frequency difference $\Delta\nu$ between the two counter-propagating laser modes. The dynamic behavior can be derived by applying the detuning either symmetrically across both beams or asymmetrically to a single beam. In both configurations, the coherent superposition of the fields is resolved using the sum-to-product identity $\cos A + \cos B = 2\cos\!\left(\tfrac{A+B}{2}\right)\cos\!\left(\tfrac{A-B}{2}\right)$.

\paragraph{Symmetric detuning.}
Consider the case in which the detuning is distributed symmetrically, shifting the optical frequencies to $\nu_1 = \nu + \Delta\nu/2$ and $\nu_2 = \nu - \Delta\nu/2$. The corresponding angular frequencies become $\omega_1 = \omega + \pi\Delta\nu$ and $\omega_2 = \omega - \pi\Delta\nu$. Evaluating the scalar electric fields near the focal region (dropping wavefront curvature and Gouy phase), the two beams are expressed as real functions:
\begin{equation}
E_1(z,t) = E_0(z,\rho) \cos(kz - \omega t - \pi\Delta\nu t),
\end{equation}
\begin{equation}
E_2(z,t) = E_0(z,\rho) \cos(-kz - \omega t + \pi\Delta\nu t) = E_0(z,\rho) \cos(kz + \omega t - \pi\Delta\nu t).
\end{equation}
Superposing the two fields, $E_{tot} = E_1 + E_2$, and mapping the arguments to $A = kz - \omega t - \pi\Delta\nu t$ and $B = kz + \omega t - \pi\Delta\nu t$, the sum-to-product identity yields
\begin{equation}
E_{tot}(z,t) = 2E_0(z,\rho) \cos\left(kz - \pi\Delta\nu t\right) \cos(-\omega t).
\end{equation}
Because the atomic motion does not follow the rapidly oscillating optical carrier, the interaction can be averaged over one optical cycle, reducing the $\cos^2(-\omega t)$ term to a constant factor of $1/2$. The effective intensity, and by extension the conservative trapping potential $U(z,\rho,t)$, is governed by the slowly varying spatial envelope. The static potential thus evolves into a moving periodic lattice:
\begin{equation}
U(z,\rho,t) = U_0 \frac{w_0^2}{w^2(z)} \exp\left(-\frac{2\rho^2}{w^2(z)}\right) \cos^2(\pi \Delta\nu t - kz).
\label{eq:moving_potential_sym}
\end{equation}
Setting $\Delta\nu=0$ recovers the stationary lattice structure.

\paragraph{Single-beam detuning.}
In laboratory implementations it is generally simpler to introduce the frequency offset into a single arm. Defining the unshifted frequency as $\nu_1 = \nu$ and the detuned frequency as $\nu_2 = \nu + \Delta\nu$ (yielding angular frequencies $\omega_1 = \omega$ and $\omega_2 = \omega + 2\pi\Delta\nu$), the interacting fields are
\begin{equation}
E_1(z,t) = E_0(z,\rho) \cos(kz - \omega t),
\end{equation}
\begin{equation}
E_2(z,t) = E_0(z,\rho) \cos(-kz - (\omega + 2\pi\Delta\nu)t) = E_0(z,\rho) \cos(kz + \omega t + 2\pi\Delta\nu t).
\end{equation}
Applying the sum-to-product identity with $A = kz - \omega t$ and $B = kz + \omega t + 2\pi\Delta\nu t$ factors the superposition:
\begin{equation}
E_{tot}(z,t) = 2E_0(z,\rho) \cos\left(kz + \pi\Delta\nu t\right) \cos\left(\omega t + \pi\Delta\nu t\right).
\end{equation}
Upon time-averaging the fast optical component $\cos^2(\omega t + \pi\Delta\nu t) \to 1/2$, the resulting effective potential is
\begin{equation}
U(z,\rho,t) = U_0 \frac{w_0^2}{w^2(z)} \exp\left(-\frac{2\rho^2}{w^2(z)}\right) \cos^2(kz + \pi\Delta\nu t).
\label{eq:moving_potential_single}
\end{equation}
This expression coincides, \CCn{up to an overall sign in the time argument of the envelope}{Il segno non e' cosmetico: inverte la DIREZIONE di trasporto. Nel caso simmetrico l'inviluppo e' $\cos^2(\pi\Delta\nu t-kz)$ e il reticolo va verso $+z$; qui e' $\cos^2(kz+\pi\Delta\nu t)$ e va verso $-z$. E la differenza non sta nel detuning simmetrico o su un solo braccio: sta in QUALE braccio porta la frequenza piu' alta. La regola generale, valida in entrambi i casi, e' che il reticolo si muove nel verso di propagazione del fascio a frequenza MAGGIORE, con $|v|=\lambda\Delta\nu/2$. Detta cosi', ``produces the same effect'' diventa precisa.}, with Eq.~\eqref{eq:moving_potential_sym}, confirming that an asymmetric detuning imposed on a single beam produces the same effect as a symmetrically distributed detuning.

\begin{figure}[h!]
\centering
\includegraphics[width=0.8\textwidth]{Cap_1_figures/frames_conveyor_belt.png}
\caption{Kinematics of the optical conveyor belt: a mutual detuning $\Delta\nu$ produces a moving standing wave.  In the moving reference frame both
beams are doppler shifted by the same amount \cite{Kuhr:03}.}
\label{fig:frames_conveyor_belt}
\end{figure}

\CCref{The macroscopic velocity of the interference pattern}{Verificato meccanicamente su tutta la tesi: in questa sezione ci sono DUE figure mai richiamate nel testo, quella dei frame del conveyor belt e quella del washboard inclinato (zero occorrenze delle loro \texttt{\textbackslash ref}). Una figura mai richiamata galleggia dove capita, e sono per giunta le due che illustrano i due punti chiave della sottosezione: il frame in cui l'onda stazionaria e' ferma, e la barriera che si abbassa sotto accelerazione. In tutto il capitolo ne ha tre: la terza e' lo spettro di assorbimento nella sezione EIT.} can be understood by considering a moving reference frame(Figure~\ref{fig:frames_conveyor_belt}). In a coordinate system translating along the $z$-axis with velocity $v$, the Doppler effect causes a red shift of the co-propagating beam and a blue shift of the counter-propagating beam. There exists a specific reference frame in which these Doppler shifts exactly compensate for the frequency difference $\Delta\nu$ imposed in the laboratory frame. In this system, the two beams have identical frequencies, and the standing wave appears stationary \cite{Kuhr:03}.

The velocity can also be derived directly in the laboratory frame by tracking a point of constant intensity, i.e.\ by setting the phase argument of the moving envelope to a constant value, $kz \pm \pi\Delta\nu t = \text{const}$. Taking the time derivative yields $k\,dz/dt \pm \pi\Delta\nu = 0$. Substituting $k = 2\pi/\lambda$ and identifying the transport velocity as $v = dz/dt$ gives
\begin{equation}
v = \frac{\lambda \Delta\nu}{2}.
\label{eq:conveyor_velocity}
\end{equation}

\subsection{Trap-Depth Scaling and Adiabatic Compression}

The dimensional properties of the trap scale dynamically when optical parameters are varied, for instance during adiabatic compression in which the beam waist is tightened towards $w_0$. Under constant optical power $P$, the trap depth obeys $U_0 \propto I_{\max} \propto P/w_0^2$. Combining this scaling with the results of Eq.~\eqref{eq:oscillation_frequencies} gives the behavior of the trapping frequencies:
\begin{equation}
\Omega_z \propto \sqrt{U_0} \propto \frac{1}{w_0}, \qquad
\Omega_{rad} \propto \frac{\sqrt{U_0}}{w_0} \propto \frac{1}{w_0^2}.
\end{equation}
The quadratic dependence of $\Omega_{rad}$, compared to the linear dependence of $\Omega_z$, induces a deformation during focusing: the anisotropy ratio $\Omega_z/\Omega_{rad}$ scales directly with $w_0$, so the radial confinement tightens faster than the axial confinement and the aspect ratio of the disc-shaped wells changes during compression.

\paragraph{Adiabatic temperature scaling.}
Provided that the variation of the trap depth remains slow compared to the inverse characteristic frequencies, the action $J \propto E/\Omega$ is an adiabatic invariant of the harmonic motion~\cite{LandauMechanics}. Modeling the mean thermal energy via the equipartition theorem ($E \propto k_B T$), this invariance implies that along each independent degree of freedom
\begin{equation}
\frac{T}{\Omega} = \text{const.}
\end{equation}
Because $\Omega_z$ and $\Omega_{rad}$ scale at different rates with $w_0$, an initially isotropic thermal distribution evolves into an anisotropic one: adiabatic compression increases the transverse temperature relative to the axial temperature.

\paragraph{Lamb--Dicke parameter.}
The axial behavior of the potential also affects the applicability of advanced laser-cooling mechanisms. The ground-state spatial width of the atomic wavepacket along the $z$-axis is
\begin{equation}
x_0 = \sqrt{\frac{\hbar}{2 m \Omega_z}},
\end{equation}
and the parameter describing the coupling of this motional state to a resonant cooling field of wavenumber $k_c = 2\pi/\lambda_c$ is the Lamb--Dicke parameter
\begin{equation}
\eta = k_c\, x_0 = k_c \sqrt{\frac{\hbar}{2 m \Omega_z}}.
\label{eq:lambdicke}
\end{equation}
The Lamb--Dicke regime provides the basis for the EIT-cooling treatment developed in Section~\ref{sec:EIT_cooling}.

\subsection{Dynamics in the Accelerated Potential}

Transporting a trapped atom between two spatial positions requires a profile of continuous acceleration followed by symmetric deceleration. This is achieved by applying a linear temporal ramp to the beat frequency $\Delta\nu(t)$.

\begin{figure}[h!]
\centering
\includegraphics[width=0.8\textwidth]{Cap_1_figures/accelerated_conveyor_belt.png}
\caption{In the frame co-moving with the accelerated lattice, the periodic potential is superimposed with a linear gradient, resulting in a tilted washboard landscape \cite{Kuhr:03}.}
\label{fig:accelerated_conveyor_belt}
\end{figure}

Transforming the Hamiltonian into the non-inertial frame of the accelerating lattice, an atom subject to an acceleration $a$ experiences a fictitious force $-ma$ (Figure~\ref{fig:accelerated_conveyor_belt}). The addition of this linear term along the optical axis modifies the effective potential:
\begin{equation}
U_{\text{eff}}(z') = U_0 \cos^2(kz') + m a z',
\label{eq:tilted_potential}
\end{equation}
where $z'$ is the coordinate in the co-moving frame and $U_0 < 0$ denotes a red-detuned trap. The linear term breaks the symmetry of the periodic profile, lowering the potential barrier in the direction of the fictitious force. Stationary minima of $U_{\text{eff}}$ exist provided the equation $dU_{\text{eff}}/dz' = 0$ admits real solutions corresponding to local minima. Differentiating yields
\begin{equation}
\frac{d U_{\text{eff}}(z')}{dz'} = -k U_0 \sin(2kz') + m a = 0 \quad \Longrightarrow \quad m a = k U_0 \sin(2kz').
\label{eq:potential_derivative}
\end{equation}
Since $|\sin(2kz')| \leq 1$, real solutions require $|ma| \leq |kU_0|$. As $|a|$ approaches this limit, the local minimum and the adjacent maximum merge into a single inflection point, causing the trapping barrier to vanish completely and the atom to escape. \CCn{The threshold acceleration for a particle with zero kinetic energy}{Giusto, ma e' un doppio best case e conviene dirlo. Primo: nel potenziale inclinato si usa $U_0$ nudo, senza inviluppo gaussiano e senza fattore radiale, quindi $a_0$ vale sull'asse e al fuoco; fuori asse la buca e' meno profonda e l'atomo sfugge prima. Secondo: a temperatura finita gli atomi superano la barriera ridotta prima che si annulli. E' il numero che limita davvero il tempo di trasporto.} is therefore:
\begin{equation}
a_0 = \frac{k |U_0|}{m}.
\label{eq:max_acceleration}
\end{equation}
This theoretical bound assumes a particle at absolute zero temperature. For a thermal cloud, the effective potential barrier $\Delta U = U_{\max} - U_{\min}$ decreases with acceleration, causing atoms with non-zero kinetic energy to escape at values significantly below $a_0$. Furthermore, the local trap frequency approaches zero ($\Omega \to 0$) as $a \to a_0$. These factors impose stringent constraints on the acceleration profiles in optical conveyor belt operations to minimize thermal loss and maintain the adiabatic confinement of the atomic wavepacket.